% !TeX program = pdflatex
\documentclass[11pt,a4paper]{article}
\usepackage{../../recursos/latex/estilo-notas}

\begin{document}
\cabecalhoaula{Extra E2}{Redução de Dimensionalidade (PCA e t-SNE)}{AME cap.~10; ISLP §12.2}

%==============================================================================
\section*{Objetivos da aula}
%==============================================================================
Ao final desta aula o aluno deve ser capaz de:
\begin{itemize}[leftmargin=1.4cm]
  \item motivar a \textbf{redução de dimensionalidade} (visualização, compressão,
        ruído, maldição da dimensionalidade);
  \item definir os \textbf{componentes principais} (PCA) como combinações lineares
        de variância máxima;
  \item resolver o PCA via \textbf{autovetores} da matriz de covariância e
        interpretar cargas e \emph{scores};
  \item escolher o número de componentes pela \textbf{proporção de variância
        explicada};
  \item reconhecer o PCA como \emph{escalonamento multidimensional} e usá-lo para
        compressão;
  \item descrever \textbf{Kernel PCA} e \textbf{t-SNE}, e os cuidados ao usar t-SNE.
\end{itemize}

\begin{emsala}
Segunda face do não supervisionado (após o agrupamento, Aula~E1). Conexão direta com
a Aula~05: a \emph{redundância} das covariáveis (dimensão intrínseca baixa) é
exatamente o que a redução de dimensionalidade explora. Pergunta de abertura:
\emph{``como resumir centenas de variáveis correlacionadas em duas ou três, sem
perder o essencial?''}
\end{emsala}

%==============================================================================
\section{Por que reduzir a dimensionalidade?}
%==============================================================================
\begin{itemize}[leftmargin=1.4cm,nosep]
  \item \textbf{Visualização}: projetar dados em 2D/3D para enxergar estrutura.
  \item \textbf{Compressão}: armazenar/transmitir com menos espaço (imagens).
  \item \textbf{Pré-processamento}: combater a \emph{maldição da dimensionalidade}
        (Aula~05) e o ruído antes de um método supervisionado.
  \item \textbf{Redundância}: muitas covariáveis correlacionadas vivem perto de uma
        subvariedade de dimensão menor --- queremos recuperá-la.
\end{itemize}

%==============================================================================
\section{Componentes principais (PCA)}
%==============================================================================
Suponha as covariáveis \emph{centradas} (média zero). O \textbf{primeiro componente
principal} é a combinação linear
\[
  Z_1 = \phi_{11}X_1+\dots+\phi_{p1}X_p = \bm{\phi}_1^\top\X
\]
de \textbf{variância máxima}, sob a restrição $\norm{\bm{\phi}_1}^2=\sum_k\phi_{k1}^2=1$
(sem ela, bastaria inflar os coeficientes). O $i$-ésimo componente $Z_i$ tem variância
máxima entre as combinações lineares unitárias \emph{não correlacionadas} com
$Z_1,\dots,Z_{i-1}$.

\begin{ideia}
Procuramos novos eixos (ortogonais) ao longo dos quais os dados \emph{mais variam}.
O 1º eixo capta o máximo de variabilidade; o 2º, o máximo do que resta sendo
ortogonal ao 1º; e assim por diante. Concentrar a variância em poucos eixos é o que
permite descartar os demais com pouca perda.
\end{ideia}

\subsection{Solução: autovetores da covariância}
Seja $\bm{C}=\bm{X}^\top\bm{X}$ (proporcional à matriz de covariância, pois os dados
têm média zero). O problema é
\[
  \max_{\bm{\phi}}\ \Var(\bm{\phi}^\top\X) = \max_{\bm{\phi}}\ \bm{\phi}^\top\bm{C}\bm{\phi}
  \quad\text{s.a.}\quad \norm{\bm{\phi}}^2=1.
\]

\begin{proposicao}[O 1º componente é o autovetor principal]\label{prop:pca}
O maximizador de $\bm{\phi}^\top\bm{C}\bm{\phi}$ sob $\norm{\bm{\phi}}=1$ é o
autovetor de $\bm{C}$ associado ao \emph{maior} autovalor, e o valor máximo da
variância é esse autovalor.
\end{proposicao}
\begin{proof}
Lagrangiano: $\mathcal{L}(\bm{\phi},\nu)=\bm{\phi}^\top\bm{C}\bm{\phi}
-\nu(\bm{\phi}^\top\bm{\phi}-1)$. Anulando o gradiente em $\bm{\phi}$:
$2\bm{C}\bm{\phi}-2\nu\bm{\phi}=0$, isto é, $\bm{C}\bm{\phi}=\nu\bm{\phi}$ ---
$\bm{\phi}$ é \emph{autovetor} de $\bm{C}$ com autovalor $\nu$. Para um autovetor
unitário, a variância é $\bm{\phi}^\top\bm{C}\bm{\phi}=\nu$. Logo o máximo é atingido
no autovetor de \emph{maior} autovalor. Os componentes seguintes saem do mesmo
argumento restrito ao complemento ortogonal, gerando os demais autovetores em ordem
decrescente de autovalor.
\end{proof}

Empilhando os autovetores (ordenados por autovalor decrescente) nas colunas de
$\bm{U}$ (a \textbf{matriz de cargas}), os \emph{scores} são
\[
  \bm{Z} = \bm{X}\bm{U}\quad (n\times p),
\]
e a $i$-ésima nova coordenada da observação $\x$ é $z_{i}=\bm{u}_i^\top\x$. Como os
autovetores são ortogonais, os componentes são \emph{não correlacionados} --- sem
informação redundante.

\subsection{Quantos componentes? Variância explicada}
Cada autovalor $\lambda_i$ é a variância capturada pelo $i$-ésimo componente. A
\textbf{proporção de variância explicada} pelos $p$ primeiros é
\begin{equation}\label{eq:pve}
  \text{PVE}(p) = \frac{\sum_{i=1}^p \lambda_i}{\sum_{i=1}^p \lambda_i}.
\end{equation}
Escolhe-se $p$ olhando o \emph{scree plot} ($\lambda_i$ vs.\ $i$, procurando o
``cotovelo'') ou exigindo PVE acima de um limiar (ex.: 90\%).

%==============================================================================
\section{Interpretação geométrica e compressão}
%==============================================================================
\subsection{PCA como escalonamento multidimensional}
O PCA também resolve: encontrar a projeção linear $\bm{Z}$ em $m<p$ dimensões que
melhor \emph{preserva as distâncias} entre pares,
\[
  \min_{\bm{Z}}\ \sum_{i,j}\big(\norm{\x_i-\x_j}^2 - \norm{\bm{z}_i-\bm{z}_j}^2\big).
\]
A solução são os $m$ primeiros componentes principais. Ou seja, PCA é a melhor
aproximação linear de baixa dimensão no sentido de distâncias --- daí ser tão útil
para visualização.

\subsection{Compressão (SVD truncada)}
Da decomposição $\bm{X}^\top\bm{X}=\bm{U}\bm{\Sigma}\bm{U}^\top$, reconstrói-se
$\bm{X}\approx \bm{Z}_p\bm{U}_p^\top$, usando só as $p$ primeiras colunas. Se os
autovalores decaem rápido, a reconstrução é quase perfeita guardando apenas
$(z_{i,1},\dots,z_{i,p})$ e as $p$ cargas $\bm{u}_1,\dots,\bm{u}_p$ --- uma
\textbf{compressão}. Aplicação clássica: comprimir imagens tratando cada canal de
cor (R, G, B) como uma matriz.

\begin{atencao}
PCA usa variância e produtos internos $\Rightarrow$ é \textbf{sensível à escala}.
\emph{Padronize} as covariáveis antes (Aula~07), senão variáveis de maior magnitude
dominam os componentes artificialmente. E PCA é \emph{não supervisionado}: a direção
de maior variância \emph{não} é necessariamente a mais útil para prever um $Y$.
\end{atencao}

%==============================================================================
\section{Além do linear: Kernel PCA e t-SNE}
%==============================================================================
\subsection{Kernel PCA}
PCA só captura estrutura \emph{linear}. Como o PCA pode ser escrito em termos de
produtos internos, aplica-se o \textbf{truque do \emph{kernel}} (Aulas~04 e~10):
substituem-se os produtos internos por um kernel $K$, fazendo PCA num espaço de
\emph{features} de alta dimensão. Resultado: componentes \textbf{não lineares}, que
capturam subvariedades curvas.

\subsection{t-SNE}
O \textbf{t-SNE} (\emph{t-distributed Stochastic Neighbor Embedding}) é uma técnica
\emph{não linear} feita para \textbf{visualização} (2D/3D). A ideia:
\begin{itemize}[leftmargin=1.4cm,nosep]
  \item converte distâncias em \textbf{probabilidades de vizinhança} no espaço
        original (kernel gaussiano) e no espaço reduzido (distribuição $t$ de
        Student, de caudas pesadas);
  \item ajusta as posições no espaço reduzido minimizando a divergência de
        \textbf{Kullback--Leibler} entre as duas distribuições;
  \item assim, \emph{preserva a estrutura local} (vizinhos próximos continuam
        próximos), revelando agrupamentos.
\end{itemize}
O principal hiperparâmetro é a \textbf{perplexidade} (ordem de grandeza do número de
vizinhos considerados).

\begin{atencao}[Cuidados com o t-SNE]
\begin{itemize}[leftmargin=1.2cm,nosep]
  \item é \textbf{só para visualização} --- não use os eixos como \emph{features}
        para um modelo supervisionado;
  \item \textbf{distâncias e tamanhos} entre grupos no gráfico \emph{não} têm
        significado --- só a vizinhança local é confiável;
  \item é \textbf{estocástico} e sensível à perplexidade: rode mais de uma vez e
        compare. Para grandes bases, UMAP é uma alternativa mais rápida.
\end{itemize}
Em resumo: PCA para \emph{compressão/pré-processamento} (linear, interpretável,
determinístico); t-SNE para \emph{enxergar} agrupamentos.
\end{atencao}

%==============================================================================
\section{Prática em Python}
%==============================================================================
\begin{lstlisting}
from sklearn.decomposition import PCA, KernelPCA
from sklearn.manifold import TSNE
from sklearn.preprocessing import StandardScaler

Xs = StandardScaler().fit_transform(X)        # padronizar e essencial

pca = PCA().fit(Xs)
print("variancia explicada acumulada:", pca.explained_variance_ratio_.cumsum())
Z = PCA(n_components=2).fit_transform(Xs)     # scores nos 2 primeiros PCs

# Kernel PCA (nao linear) e t-SNE (visualizacao)
Zk  = KernelPCA(n_components=2, kernel="rbf", gamma=0.1).fit_transform(Xs)
Zt  = TSNE(n_components=2, perplexity=30, init="pca",   # init explicito:
           random_state=0).fit_transform(Xs)            # o padrao ja mudou uma vez
\end{lstlisting}
Os notebooks \texttt{Aula prática E2} e \texttt{Exemplo - PCA} ilustram os
componentes principais e comparam com o
t-SNE na mesma base.

%==============================================================================
\section*{Resumo}
%==============================================================================
\begin{itemize}[leftmargin=1.4cm]
  \item Reduzir dimensão: visualizar, comprimir, remover ruído, explorar
        \emph{redundância} (Aula~05).
  \item \textbf{PCA}: componentes $=$ combinações lineares de variância máxima
        $=$ autovetores da covariância (Prop.~\ref{prop:pca}); \emph{scores}
        $\bm{Z}=\bm{X}\bm{U}$, não correlacionados.
  \item Número de componentes pela \textbf{PVE} \eqref{eq:pve} / \emph{scree plot}.
  \item PCA $=$ melhor projeção linear que preserva distâncias (MDS); base da
        \textbf{compressão} por SVD truncada. \emph{Padronize}.
  \item \textbf{Kernel PCA}: PCA não linear via \emph{kernel}. \textbf{t-SNE}:
        visualização que preserva vizinhança local --- com cuidados (estocástico,
        só visual, distâncias entre grupos sem sentido).
\end{itemize}

\paragraph{Leitura recomendada.} [AME] Capítulo~10 (§10.1 PCA e MDS, §10.1.2
compressão de imagens, §10.2 Kernel PCA); [ISLP] §12.2 (\emph{Principal Components
Analysis}).

\end{document}
